% FixedCase_D347_k1_Proof.tex
% MathWise DI^2 — Fixed Case Only - Erdős 124 
% D = {3,4,7}, k = 1
% No general BEGL result claimed.

\documentclass[11pt]{article}

\usepackage{amsmath,amssymb,amsthm}
\usepackage{geometry}
\usepackage{array}
\usepackage{booktabs}
\usepackage{longtable}

\geometry{margin=1in}

\newtheorem{theorem}{Theorem}
\newtheorem{lemma}{Lemma}
\newtheorem{proposition}{Proposition}
\newtheorem{certificate}{Certificate}

\newcommand{\SubSum}{\operatorname{SubSum}}
\newcommand{\ord}{\operatorname{ord}}

\title{Fixed-Case Certificate for \(D=\{3,4,7\},\ k=1\)}
\author{MathWise DI\(^2\) Fixed-Case Proof Package}
\date{}

\begin{document}

\maketitle

\section{Scope and Non-Claims}

This proof package concerns only the fixed case
\[
D=\{3,4,7\},\qquad k=1.
\]

The result proved here is only
\[
[582,\infty)\subset A_1(\{3,4,7\}).
\]

No general BEGL conjecture is claimed.

No result is claimed for any other set \(D\).

No result is claimed for \(k\ne 1\).

No lower endpoint below \(582\) is claimed.

No optimality claim is made.

No density or asymptotic claim is made.

\section{Theorem Statement}

Let \(A_1(\{3,4,7\})\) denote the set of integers representable as a subset sum of distinct powers generated by
\[
3^a,\qquad 4^b,\qquad 7^c,
\]
with positive exponents.

\begin{theorem}[Fixed case \(D=\{3,4,7\}, k=1\)]
\[
[582,\infty) \subset A_1(\{3,4,7\}).
\]
\end{theorem}

\section{Initial Power List and Interval Certificate}

The initial ordered power list is
\[
b_1,\ldots,b_{15}
=
3,4,7,9,16,27,49,64,81,243,256,343,729,1024,2187.
\]

The certified interval is
\[
[582,4460]\subset \SubSum(b_1,\ldots,b_{15}).
\]

The fixed certificate parameters are
\[
M=582,
\]
\[
L_{15}=4460-582=3878,
\]
\[
S_0=5042,
\]
\[
b_{15}=2187,
\]
\[
b_{16}=2401.
\]

Thus the proof begins from the finite interval certificate
\[
[582,4460]\subset \SubSum(b_1,\ldots,b_{15}).
\]

\section{Direct Interval-Overlap Lemma}

\begin{lemma}[Interval-overlap propagation]
Let
\[
[M,N]\subset \SubSum(S)
\]
for a finite set of distinct allowed terms \(S\). Let \(p\) be a new unused allowed term. If
\[
p\le N-M+1,
\]
then
\[
[M,N+p]\subset \SubSum(S\cup\{p\}).
\]
\end{lemma}

\begin{proof}
The interval \([M,N]\) is already covered. Adding \(p\) to each representation covers
\[
[M+p,N+p].
\]
If
\[
p\le N-M+1,
\]
then
\[
M+p\le N+1,
\]
so the intervals \([M,N]\) and \([M+p,N+p]\) overlap or touch. Therefore their union covers
\[
[M,N+p].
\]
Hence
\[
[M,N+p]\subset \SubSum(S\cup\{p\}).
\]
\end{proof}

Let \(B\) be the ordered set of powers generated by \(3,4,7\):
\[
B=\{3^a:a\ge1\}\cup\{4^b:b\ge1\}\cup\{7^c:c\ge1\}.
\]

For \(t\in B\), define
\[
S(t)=\sum_{\substack{p\in B\\p<t}}p.
\]

For a future ordered power \(t>2187\), the direct interval-overlap condition is
\[
S(t)-t\ge1163.
\]

Equivalently, if \(b_j=t\) is the next future term, then
\[
b_j\le L_{j-1}+1.
\]

\section{Weighted Next-Gap Identities}

For any future power \(t>2187\), define the next-power gaps
\[
\Delta_3(t)=\min\{3^a:3^a>t\}-t,
\]
\[
\Delta_4(t)=\min\{4^b:4^b>t\}-t,
\]
\[
\Delta_7(t)=\min\{7^c:7^c>t\}-t.
\]

Only the two competing bases are used in each case.

\begin{lemma}[Weighted next-gap identities]
The following identities hold.

If \(t=3^r\), \(r\ge8\), then
\[
S(t)-t-1163
=
\frac{2\Delta_4+\Delta_7-7002}{6}.
\]

If \(t=4^s\), \(s\ge6\), then
\[
S(t)-t-1163
=
\frac{3\Delta_3+\Delta_7-7002}{6}.
\]

If \(t=7^q\), \(q\ge4\), then
\[
S(t)-t-1163
=
\frac{3\Delta_3+2\Delta_4-7002}{6}.
\]
\end{lemma}

\begin{proof}
For \(t=3^r\), let
\[
4^a=3^r+\Delta_4,\qquad 7^b=3^r+\Delta_7
\]
be the least larger \(4\)- and \(7\)-powers. Then
\[
S(3^r)
=
\frac{3^r-3}{2}
+
\frac{4^a-4}{3}
+
\frac{7^b-7}{6}.
\]
Substituting \(4^a=3^r+\Delta_4\) and \(7^b=3^r+\Delta_7\) gives
\[
S(3^r)-3^r-1163
=
\frac{2\Delta_4+\Delta_7-7002}{6}.
\]

The cases \(t=4^s\) and \(t=7^q\) are identical algebraic substitutions. For \(t=4^s\), use
\[
3^a=4^s+\Delta_3,\qquad 7^b=4^s+\Delta_7.
\]
For \(t=7^q\), use
\[
3^a=7^q+\Delta_3,\qquad 4^b=7^q+\Delta_4.
\]
This yields the stated identities.
\end{proof}

\section{Weighted Gap Lemma}

\begin{lemma}[Weighted Gap Lemma]
For every future power \(t>2187\) generated by \(3,4,7\), the corresponding weighted next-power gap satisfies:

If \(t=3^r\), \(r\ge8\), then
\[
2\Delta_4+\Delta_7\ge7002.
\]

If \(t=4^s\), \(s\ge6\), then
\[
3\Delta_3+\Delta_7\ge7002.
\]

If \(t=7^q\), \(q\ge4\), then
\[
3\Delta_3+2\Delta_4\ge7002.
\]
\end{lemma}

\begin{proof}
The proof is by finite full paired modular certificates. The full paired modular certificates are supplied in
\texttt{WEIGHTED\_GAP\_CERTIFICATE\_TABLES.md} and replayed by
\texttt{verify\_weighted\_gap\_certificates.py}.

A failure of the first inequality gives a bounded full system
\[
4^a-3^r=u,
\]
\[
7^b-3^r=v,
\]
\[
2u+v\le7001,
\]
\[
r\ge8,
\]
with
\[
1\le u\le3500,\qquad 1\le v\le7000.
\]

A failure of the second inequality gives
\[
3^a-4^s=u,
\]
\[
7^b-4^s=v,
\]
\[
3u+v\le7001,
\]
\[
s\ge6,
\]
with
\[
1\le u\le2333,\qquad 1\le v\le7000.
\]

A failure of the third inequality gives
\[
3^a-7^q=u,
\]
\[
4^b-7^q=v,
\]
\[
3u+2v\le7001,
\]
\[
q\ge4,
\]
with
\[
1\le u\le2333,\qquad 1\le v\le3500.
\]

The stabilized residue filters are as follows.

For Case A, \(r\ge8\), so
\[
3^r\equiv0\pmod{6561}.
\]
Hence
\[
4^a\equiv u\pmod{6561},
\]
\[
7^b\equiv v\pmod{6561}.
\]
Use
\[
\ord_{6561}(4)=2187,\qquad \ord_{6561}(7)=2187.
\]

For Case B, \(s\ge6\), so
\[
4^s\equiv0\pmod{4096}.
\]
Hence
\[
3^a\equiv u\pmod{4096},
\]
\[
7^b\equiv v\pmod{4096}.
\]
Use
\[
\ord_{4096}(3)=1024,\qquad \ord_{4096}(7)=512.
\]

For Case C, \(q\ge4\), so
\[
7^q\equiv0\pmod{2401}.
\]
Hence
\[
3^a\equiv u\pmod{2401},
\]
\[
4^b\equiv v\pmod{2401}.
\]
Use
\[
\ord_{2401}(3)=2058,\qquad \ord_{2401}(4)=1029.
\]

The attached replay verifier reconstructs the residue-admissible systems and confirms
\[
\text{Case A remaining}=0,\qquad
\text{Case B remaining}=0,\qquad
\text{Case C remaining}=0.
\]
The certificate checks the full paired systems, not only the difference equation. Thus every possible failure system is excluded, and the three weighted inequalities hold.
\end{proof}

\section{Full Paired Modular Certificates}

The certificate method eliminates full paired systems. The compressed certificate table and replay verifier are external-review attachments to this proof package.

\subsection*{Case A}

The full system is
\[
4^a-3^r=u,
\]
\[
7^b-3^r=v,
\]
\[
r\ge8.
\]

After the stabilized residue filter modulo \(6561\), there are
\[
1{,}361{,}889
\]
residue-admissible systems.

They are eliminated by full paired modular certificates using the following moduli:

\[
5,\ 13,\ 19,\ 31,\ 37,\ 41,\ 43,\ 73,\ 109,\ 487,\ 2917,\ 17497.
\]

The final Case A remaining row before the expanded modulus is
\[
(u,v)=(1909,214),
\]
with forced classes
\[
a\equiv1464\pmod{2187},
\]
\[
b\equiv733\pmod{2187}.
\]

It is eliminated modulo
\[
17497.
\]

\subsection*{Case B}

The full system is
\[
3^a-4^s=u,
\]
\[
7^b-4^s=v,
\]
\[
s\ge6.
\]

After the stabilized residue filter modulo \(4096\), there are
\[
255{,}938
\]
residue-admissible systems.

They are eliminated by full paired modular certificates using the following moduli:

\[
5,\ 11,\ 13,\ 17,\ 19,\ 29,\ 31,\ 37,\ 41,\ 43,\ 257.
\]

\subsection*{Case C}

The full system is
\[
3^a-7^q=u,
\]
\[
4^b-7^q=v,
\]
\[
q\ge4.
\]

After the stabilized residue filter modulo \(2401\), there are
\[
1{,}501{,}666
\]
residue-admissible systems.

They are eliminated by full paired modular certificates using the following moduli:

\[
5,\ 13,\ 17,\ 19,\ 29,\ 31,\ 37,\ 41,\ 43,\ 73,\ 197,\ 1373.
\]

The final two Case C remaining rows before modulo \(1373\) are
\[
(u,v)=(386,681),
\]
\[
(u,v)=(1844,681).
\]

Both are eliminated modulo
\[
1373.
\]

\subsection*{Least-Larger Audit}

The full paired modular certificates eliminate stronger systems than the least-larger next-power systems. Therefore the least-larger inequalities are not needed in the final modular proof.

\section{Direct Propagation}

By the Weighted Gap Lemma and the weighted next-gap identities, for every future ordered power
\[
t>2187
\]
generated by \(3,4,7\),
\[
S(t)-t\ge1163.
\]

Equivalently, for every future ordered power \(b_j\),
\[
b_j\le L_{j-1}+1.
\]

Starting from
\[
[582,4460]\subset \SubSum(b_1,\ldots,b_{15}),
\]
the interval-overlap lemma applies successively to every later ordered power.

Therefore the covered interval propagates indefinitely:
\[
[582,\infty)\subset A_1(\{3,4,7\}).
\]

\section{Final Conclusion}

Thus, for the fixed case
\[
D=\{3,4,7\},\qquad k=1,
\]
we conclude
\[
[582,\infty) \subset A_1(\{3,4,7\}).
\]

No general BEGL result is claimed.

No result for other \(D\) is claimed.

No result for \(k\ne1\) is claimed.

No endpoint below \(582\) is claimed.

\section{Appendix: Survivor-System Compatibility}

The survivor-system framework was separately closed during the fixed-case analysis.

The recorded survivor-system closure status is:

\[
\text{Case 1: closed},
\]
\[
\text{Case 2: closed},
\]
\[
\text{Case 3: closed}.
\]

The augmented coarse candidate closure status is:

\[
\text{Case 1: }51/51\text{ closed},
\]
\[
\text{Case 2: }52/52\text{ closed},
\]
\[
\text{Case 3: }38/38\text{ closed}.
\]

This survivor-system closure remains compatible with the fixed-case result.

However, the final proof above does not require the survivor-system framework, because the direct interval-overlap route supplies an independent certificate path:

\[
\text{finite interval certificate}
+
\text{interval-overlap lemma}
+
\text{Weighted Gap Lemma}
\Longrightarrow
[582,\infty)\subset A_1(\{3,4,7\}).
\]

\end{document}
